\documentclass[11pt,a4paper]{article}
\usepackage[margin=25mm]{geometry}
\usepackage[T1]{fontenc}
\usepackage{amsmath,amssymb,amsthm,mathtools,microtype,booktabs}
\usepackage[dvipsnames]{xcolor}
\usepackage[colorlinks=true,linkcolor=MidnightBlue,urlcolor=MidnightBlue,citecolor=MidnightBlue]{hyperref}
\hypersetup{pdftitle={From Euler's restricted-prime sums to Eisenstein norms},pdfauthor={Roy van Rijn}}
\usepackage{enumitem,verbatim}
\setlist{nosep}
\newtheorem{theorem}{Theorem}[section]
\newtheorem{lemma}[theorem]{Lemma}
\newtheorem{proposition}[theorem]{Proposition}
\newtheorem{conjecture}[theorem]{Conjecture}
\theoremstyle{remark}\newtheorem{remark}[theorem]{Remark}
\newcommand{\Z}{\mathbb Z}\newcommand{\Q}{\mathbb Q}\newcommand{\F}{\mathbb F}\newcommand{\C}{\mathbb C}
\newcommand{\OO}{\mathcal O}\renewcommand{\AA}{\mathcal A}\newcommand{\PP}{\mathcal P}
\newcommand{\NN}{\mathcal N}
\newcommand{\nrd}{\operatorname{nrd}}\newcommand{\PGL}{\operatorname{PGL}}
\title{From Euler's restricted-prime sums to Eisenstein norms}
\author{Roy van Rijn}
\date{17 November 2026}
\begin{document}
\maketitle

\begin{abstract}
We start from the restricted-prime sum theorem of Howe, Huffman,
Kohlmeyer and Kruer and follow its square-lattice argument into the
Eisenstein lattice. Let $\AA$ contain $1$ and the positive integers whose
prime factors are all $1$ modulo $3$. The first answer is that
$\AA+\AA$ contains exactly the positive integers $2$ modulo $6$.
The ramified prime $3$ then leads to a mixed problem: $\AA+3\!\cdot\!\AA$
is exactly the positive class $4$ modulo $6$. Quaternionic norm counts,
an odd-component filter and a six-point calculation at $3$ prove the
large-target statements; exact finite certificates complete them.
Together, the two identities give the binary and higher sumsets of
primitive Eisenstein norm values. Four such values suffice for every
positive integer, while seven summands from $\AA$ suffice; both bounds
are sharp. We then extract a reusable counting criterion with explicit
hypotheses for an effective cutoff, and end with the balanced-summand
question that remains open.
\end{abstract}

\section{The Eisenstein split-prime cores}
Let $Q(x,y)=x^2-xy+y^2$, and put
\[
 \AA=\{n\ge1:p\mid n\Longrightarrow p\equiv1\pmod3\}.
\]
Every member of $\AA$ is $1\bmod6$.  The ramified dilation
$3\!\cdot\!\AA=\{3a:a\in\AA\}$ has residue $3\bmod6$.

\begin{theorem}\label{thm:main}
One has
\[
 \AA+\AA=\{n>0:n\equiv2\pmod6\},\qquad
 \AA+(3\!\cdot\!\AA)=\{n>0:n\equiv4\pmod6\}.
\]
\end{theorem}

\paragraph{A reduction that helps both problems.}
Put $\PP=\{p\text{ prime}:p\equiv2\pmod3\}$. For either admissible residue
class, write $n=gn_0$, where $g\in\AA$ contains every prime-power factor
of $n$ with prime $1$ modulo $3$. Since $g\equiv1\pmod6$, the target
$n_0$ has the same residue as $n$. A decomposition of either type for
$n_0$ scales by $g$ to the corresponding decomposition of $n$.
We can therefore concentrate on
\begin{equation}\label{eq:reduced}
n\equiv2\text{ or }4\pmod6,\qquad
p\mid n\Longrightarrow p\in\PP.
\end{equation}
This reduction will matter in the divisor estimates and in the exact
calculation at $3$. The finite certificates cover all targets in their
respective ranges, so they need no reduction.

\section{Counting with Eisenstein norms}\label{sec:norm}
We begin by finding plenty of candidate pairs. A quaternion packages two
Eisenstein norms into one multiplicative norm, which lets us count the
representations exactly. This count is the supply of pairs from which we
will remove the bad ones.

For positive integers, write $\sigma(n)=\sum_{d\mid n}d$ and
$\tau(n)=\sum_{d\mid n}1$. A term $\sigma(n/4)$ is interpreted as zero
when $4\nmid n$. Let $\omega^2+\omega+1=0$, $R=\Z[\omega]$, and
\[
Q(x,y)=x^2-xy+y^2=N_{\Q(\omega)/\Q}(x+y\omega).
\]
For instance, $Q(3,1)=7$ and $Q(4,1)=13$, so $20$ is represented by
$Q(x,y)+Q(u,v)$ with both norms already in $\AA$.

Consider the definite quaternion algebra and order
\[
B=\Q(\omega)\oplus\Q(\omega)j,\qquad
j^2=-1,\quad jz=\bar z j,\qquad \OO=R\oplus Rj.
\]
Its reduced norm is
\[
\nrd(z+w j)=N(z)+N(w).
\]
In the basis $1,\omega,j,\omega j$, the trace pairing has determinant $-9$.
Consequently $\OO$ has reduced discriminant $3$. Definiteness and the parity
of the ramification set imply that the algebra discriminant is $3$, so
$\OO$ is maximal.

Every complex number is within squared distance $1/3$ of an Eisenstein
integer: the equilateral triangles in the lattice have circumradius
$1/\sqrt3$. Rounding the two complex coordinates independently gives
\[
\forall x\in B\otimes\mathbb R\quad\exists q\in\OO:
\quad\nrd(x-q)\le2/3<1.
\]
The one-coordinate rounding bound also makes $R$ norm Euclidean, hence
a unique factorization domain. The two-coordinate bound makes $\OO$
norm Euclidean on either side. In a nonzero integral right
ideal, choose $\alpha$ of least positive norm and divide each $\beta$ by
$\alpha$ on the left. A nonzero remainder would have smaller norm, so the
ideal is $\alpha\OO$. Fractional ideals reduce to this case by scaling;
the right ideal class number is one.

The units are precisely $R^\times\cup R^\times j$, and there are twelve.
For $3\nmid n$ let
\[
\mathcal Q_n=\{\alpha\in\OO:\nrd(\alpha)=n\},\qquad r(n)=|\mathcal Q_n|.
\]
At a prime $\ell\mid n$ the order is locally $M_2(\Z_\ell)$. There are
$1+\ell+\cdots+\ell^e$ integral right ideals of reduced norm $\ell^e$,
by the index-$\ell^e$ sublattice count in $\Z_\ell^2$. Multiplication of
the local counts and principality give $\sigma(n)$ right ideals of norm
$n$, each with twelve generators of that norm. Hence
\begin{equation}\label{eq:count}
r(n)=12\sigma(n)\qquad(3\nmid n).
\end{equation}
This is the usual all-integral-ideal normalization of a Brandt operator;
see \cite[Proposition 41.3.1]{Voight}.

For $n\equiv2\pmod3$, both component norms in every representation
$n=Q(x,y)+Q(u,v)$ are $1\pmod3$, because
$Q(x,y)\equiv(x+y)^2\pmod3$ is either $0$ or $1$. In particular both
summands are positive and neither is divisible by $3$.

\section{Keeping the unwanted primes out}\label{sec:spectral}
A norm representation is only a candidate. A prime $p\equiv2\pmod3$
can occur to even exponent in an Eisenstein norm, and that already keeps
the value out of $\AA$. The useful feature is that divisibility of the
norm forces divisibility of both coordinates.

For $p\in\PP$, the quotient $R/pR$ is a field, including when $p=2$.
Therefore
\begin{equation}\label{eq:anisotropic}
p\mid Q(x,y)\quad\Longleftrightarrow\quad p\mid x\text{ and }p\mid y.
\end{equation}
We first establish the unfiltered Hecke estimate. Section~\ref{sec:odd-ramified}
will impose oddness and derive the local bounds used in the sieve.

\subsection{The finite quotient and its Hecke operator}
Fix an odd $p\in\PP$, so $p\ge5$, and first suppose $p\nmid n$.
Reduction gives $\OO/p\OO\simeq M_2(\F_p)$, with norm equal to determinant.
Write
\[
G_p=\PGL_2(\F_p),\quad \Gamma_p=\rho_p(\OO^\times),\quad
Y_p=\Gamma_p\backslash G_p.
\]
Only $\pm1$ among the units reduce to scalars, so
$|\Gamma_p|=6$ and $|Y_p|=p(p^2-1)/6$. Equip $\C[Y_p]$ with counting
measure and define
\begin{equation}\label{eq:hecke}
(T_nf)(\Gamma_p g)=\frac1{12}\sum_{\alpha\in\mathcal Q_n}
f(\Gamma_p\rho_p(\alpha)g).
\end{equation}
Right multiplication by a unit permutes $\mathcal Q_n$, proving that this
is well-defined. It commutes with the right $G_p$ action. Conjugation
preserves $\mathcal Q_n$ and is projective inversion, so $T_n$ is
self-adjoint. Also $T_n1=\sigma(n)1$.

The determinant character
$\chi_p(\Gamma_p g)=(\det(g)/p)$ is well-defined, because the units have
norm one and scalar matrices have square determinant. Directly,
$T_n\chi_p=(n/p)\sigma(n)\chi_p$.

\subsection{The spectral bound}
This is the main imported ingredient. It turns the large supply of norm
representations into control of their distribution modulo a prime. The
normalization matters, so we identify the concrete operator before using
Jacquet--Langlands and the weight-two Ramanujan bound.

There are three things to check: the finite operator must be the stated
Hecke correspondence, the two large character components must be removed,
and the remaining spectrum must transfer to weight two. The proof below
checks these points in that order. Readers following the arithmetic first
can use the bound and continue with the normalizer average in the next
subsection.

\begin{lemma}\label{lem:spectral}
On $\{1,\chi_p\}^\perp\subset\C[Y_p]$, the operator $T_n$ has norm at most
$\tau(n)\sqrt n$, provided $(n,3p)=1$.
\end{lemma}
\begin{proof}
Class number one identifies the adelic quotient at the reduction kernel
with $\Gamma_p\backslash G_p$.  Indeed an adelic ideal principalizes;
and a rational projective element in all local order-unit groups can
be rescaled by a rational scalar to an order unit, since the local
valuations of its norm are even.  Its invariant measure is counting
measure up to normalization.

The all-integral-ideal correspondence divides the principal shell by
twelve units.  Principalization acts by $\rho_p(\alpha^{-1})$;
projectively this equals $\rho_p(\bar\alpha)$, and conjugation
permutes the shell.  Thus its operator is exactly \eqref{eq:hecke},
including the prime-power normalization.  Every global norm character
restricts to a character of $G_p$.  A character of
$\mathrm{GL}_2(\F_p)$ kills elementary unipotents and hence factors
through determinant; descent to $\mathrm{PGL}_2$ leaves only the
trivial and quadratic determinant characters.  Both have been
identified above.  Removing them, Lemma~\ref{lem:hecke-transfer}
therefore applies for $(n,3p)=1$ and gives the asserted bound.
\end{proof}

\subsection{Torus-normalizer averaging}
The image $C_p=(R/pR)^\times/\F_p^\times$ is a nonsplit torus. Its
normalizer is $K_p=C_p\cup C_p\rho_p(j)$, of order $2(p+1)$.
All twelve order units lie in $R^\times\cup R^\times j$, so
$\Gamma_p\subset K_p$. Put
\[
O_p=\Gamma_p\backslash K_p\subset Y_p,\qquad f_p=1_{O_p},\qquad
\mu_p=\frac{|O_p|}{|Y_p|}=\frac2{p(p-1)}.
\]
An invertible quaternion has one of its two Eisenstein components zero
modulo $p$ exactly when its projective reduction lies in $K_p$.
Let $b_p(n)$ count the unfiltered representations spoiled by $p$. Consequently
\[
b_p(n)=12(T_nf_p)(\Gamma_p).
\]
Both $f_p$ and $T_nf_p$ are right-$K_p$ invariant. Therefore
\[
(T_nf_p)(\Gamma_p)=\frac{\langle f_p,T_nf_p\rangle}{|O_p|}.
\]
The norm map $\F_{p^2}^\times\to\F_p^\times$ is onto, so $O_p$ has
equally many determinant-square and determinant-nonsquare points.
Thus $f_p-\mu_p1$ is orthogonal to both $1$ and $\chi_p$, and has
squared norm $|O_p|(1-\mu_p)$. The spectral lemma gives
\begin{equation}\label{eq:local}
\frac{b_p(n)}{r(n)}
\le\frac2{p(p-1)}+
\frac{\tau(n)\sqrt n}{\sigma(n)}
\left(1-\frac2{p(p-1)}\right).
\end{equation}

\section{An odd-component filter and the mixed case}\label{sec:odd-ramified}
The prime $2$ need not take a share of the error budget. We can remove
its obstruction exactly by keeping only odd norm components. This helps
both binary problems and leaves room for the extra condition at $3$ in
the mixed case.

Throughout this section, $n$ is even and $3\nmid n$. Write $n=2^e m$
with $e\ge1$ and $m$ odd, and put
\[
S(n)=\sigma(n)-\sigma(n/4),
\]
where a term with nonintegral argument is zero. Let $\mathcal Q_n^o$
consist of the representations in $\mathcal Q_n$ for which both
Eisenstein component norms are odd.

\subsection{Selecting odd components}
Since $2$ is inert in $\Z[\omega]$, an even component norm has both
coordinates even. The two component norms of an even target have the
same parity. Thus the representations excluded from $\mathcal Q_n^o$
are exactly the scalar multiples $2\alpha$, with
$\alpha\in\mathcal Q_{n/4}$. Consequently
\begin{equation}\label{eq:odd-count}
R_o(n):=|\mathcal Q_n^o|=12S(n),\qquad
S(n)=3\cdot2^{e-1}\sigma(m)\ge\frac32n.
\end{equation}
Both retained norms are positive, since zero is even.

For an odd inert prime $p$, let $b_p^o(n)$ count the retained
representations in which $p$ divides either component norm. When $p\mid n$,
all four coordinates are divisible by $p$ whenever one component norm is.
Scaling by $p$ preserves oddness. Therefore
\[
b_p^o(n)=0\quad(p^2\nmid n),\qquad
b_p^o(n)=12S(n/p^2)\quad(p^2\mid n),
\]
and the usual local divisor-sum calculation gives
\begin{equation}\label{eq:odd-dividing}
\frac{b_p^o(n)}{R_o(n)}<\frac1{p^2}.
\end{equation}

\subsection{The odd-component spectral estimate}
Fix an odd inert prime $p\nmid n$, and use the quotient $Y_p$ and
operators $T_u$ of Section~\ref{sec:spectral}. At this odd prime, the
scalar $2$ reduces projectively to the identity. Thus
\[
D_n:=T_n-T_{n/4}
\]
is exactly the operator obtained by retaining only $\mathcal Q_n^o$;
again the second term is zero when $4\nmid n$. Its constant eigenvalue
is $S(n)$, and its determinant-character eigenvalue is
$(n/p)S(n)$, since $4$ is a square modulo $p$.

On the orthogonal complement of these two characters, consider a
simultaneous cuspidal eigenvector. At $2$, write the normalized Satake
roots as $u,v$, with $|u|=|v|=1$ and $uv=1$. The last identity follows
from the trivial central character of the projective quotient. The
all-ideal eigenvalue at $2^e$, minus the eigenvalue at $2^{e-2}$, is
\[
2^{e/2}\left(u^e+v^e+
\frac12\sum_{s=1}^{e-1}u^s v^{e-s}\right).
\]
For $e=1$ the sum is empty. To obtain this expression, use $uv=1$ to
identify the terms from $2^{e-2}$ with the interior terms from $2^e$;
the former have an additional factor $1/2$. Its absolute value is at most $((e+3)/2)2^{e/2}$. Multiplicativity and
Lemma~\ref{lem:spectral} at the odd factors of $m$ now give
\begin{equation}\label{eq:odd-operator}
\|D_n\|_{\{1,\chi_p\}^\perp}
\le\frac{e+3}{2}\tau(m)\sqrt n.
\end{equation}
This estimate uses cancellation before taking absolute values.

Suppose further that $n$ is reduced, so every prime factor lies in
$\PP$. The ratio of the right side of \eqref{eq:odd-operator} to $S(n)$ is
\[
n^{-3/10} H(e)\prod_{q^a\Vert m}F_q(a),\qquad
H(e)=\frac{e+3}{3\,2^{e/5}},\qquad
F_q(a)=\frac{(a+1)q^{4a/5}}{1+q+\cdots+q^a}.
\]
The ratio $H(e+1)/H(e)$ decreases with $e$, and exact fifth-power
comparisons show that $H(e)$ attains its maximum for $e\ge1$ at $e=4$.
For $q\ge32$, one has $F_q(a)\le(a+1)q^{-a/5}\le1$.
For the remaining odd inert primes, successive ratios are bounded by
$((a+2)/(a+1))q^{-1/5}$. Once this upper bound is at most one it remains
so. Checking the preceding exponents gives
\[
\begin{array}{c|rrrrr}
q&5&11&17&23&29\\
\text{maximizing exponent}&2&1&1&1&0.
\end{array}
\]
The product of the fifth powers of these maxima, including $H(4)^5$, is
\[
\frac{2246611641466280158984375}{52336356030243694706688}
<\frac{59049}{1024}=\left(\frac94\right)^5.
\]
Thus the ratio in question is less than $(9/4)n^{-3/10}$.

The normalizer-orbit averaging argument of Section~\ref{sec:spectral}
applies to $D_n$ with degree $S(n)$. Its orbit indicator has the same
orthogonality to both character constituents. Therefore
\begin{equation}\label{eq:odd-local}
\frac{b_p^o(n)}{R_o(n)}
\le\frac{2}{p(p-1)}+\frac94n^{-3/10}.
\end{equation}
When $p\mid n$, \eqref{eq:odd-dividing} is already smaller than the first
term, so \eqref{eq:odd-local} holds for every odd inert prime.

\subsection{A six-point quotient at the ramified prime}
For $\AA+\AA$, a target $2$ modulo $3$ forces both component norms
to be $1$ modulo $3$. The mixed target is $1$ modulo $3$, so exactly one
component norm is divisible by $3$. To make it $3b$ with $b\in\AA$,
we must keep its exponent of $3$ equal to one. We are left with a concrete
task: count the representations with a component norm divisible by $9$.
A six-point quotient gives that count exactly, without another spectral
estimate. Put
\[
\overline R=R/3R,\qquad
G_3=(\OO/3\OO)^\times/\F_3^\times,\qquad
\Gamma_3=\operatorname{image}(\OO^\times),\qquad
Y_3=\Gamma_3\backslash G_3.
\]
For $z=x+y\omega$ one has $N(z)\equiv(x+y)^2\pmod3$. Among the $81$
residue quaternions $z+w j$, exactly nine have norm zero: both of the
corresponding coordinate sums vanish. The other $72$ are units, since
quaternionic conjugation supplies their inverses after dividing by norm.
Hence $|G_3|=36$. Only $\pm1$ of the twelve order units become scalars,
so $|\Gamma_3|=6$ and $|Y_3|=6$.

The norm modulo $3$ is well-defined projectively and is one on
$\Gamma_3$. Each norm class $Y_3^+$ and $Y_3^-$ has three points.
Indeed, norm one occurs when exactly one of the two coordinate sums is
nonzero, and norm two when both are nonzero; each condition describes
$36$ unit quaternions before projectivizing.

The six units of $\overline R$ are exactly the reductions of
$\pm1,\pm\omega,\pm\omega^2$. Every point of $Y_3^-$ may therefore be
represented by $1+t j$, with $t\in\overline R^\times$: first normalize
its first component by a left Eisenstein unit. Left multiplication by
$j$, followed by the same normalization, sends $t$ to $-t$, because
$t\bar t=1$. Left multiplication by an Eisenstein unit leaves the
normalized ratio unchanged. These exhaust all order units, so the three
left cosets are exactly
\[
\Gamma_3(1+j),\qquad
\Gamma_3(1+\omega j),\qquad
\Gamma_3(1+\omega^2j).
\]
Right multiplication by $\omega$, followed by left normalization, sends
$t$ to $t\omega$ and cycles these three points. Thus the right
$\Gamma_3$-action is transitive on $Y_3^-$.

The same ideal-class argument as in Lemma~\ref{lem:spectral} realizes
$Y_3$ as an adelic quotient, now taking the kernel of reduction
$\OO_3^\times\to G_3$ at $3$ and maximal-order units elsewhere.
For $3\nmid u$, its all-ideal Hecke operator is still
\[
(T_uf)(\Gamma_3 g)=\frac1{12}\sum_{\alpha\in\mathcal Q_u}
 f(\Gamma_3\rho_3(\alpha)g),
\]
with degree $\sigma(u)$, counting measure, self-adjointness and the
prime-power recurrence already established.

For an inert prime $\ell$, $T_\ell$ interchanges the two norm classes.
Its value on the base-point indicator is constant on the opposite
three-point class, by the transitive right $\Gamma_3$-action just proved.
Self-adjointness and degree $\ell+1$ make that value $(\ell+1)/3$.
Equivariance under $G_3$ transports this conclusion to every point. Thus,
with the two norm classes ordered first,
\[
T_\ell=\frac{\ell+1}{3}
\begin{pmatrix}0&J_3\\J_3&0\end{pmatrix},
\]
where $J_3$ is the $3$ by $3$ all-ones matrix.
It has eigenvalues $\ell+1$ and $-(\ell+1)$ on the constant and norm-sign
characters, and acts as zero on the four-dimensional space $W$ of
functions with sum zero on each norm class.

The all-ideal recurrence shows that the scalar on $W$ at $\ell^a$ is
zero for odd $a$ and $(-\ell)^{a/2}$ for even $a$. For a reduced target
$n$, multiplicativity gives
\[
\lambda_n=
\begin{cases}
0,&n\text{ is not a square},\\
(-1)^{\Omega(n)/2}\sqrt n,&n\text{ is a square},
\end{cases}
\]
where $\Omega$ counts prime factors with multiplicity.

\subsection{The exact count and its odd-component bound}
For the unfiltered identity, allow any positive reduced target
$n\equiv1\pmod3$, without an evenness requirement. If
$9\mid Q(x,y)$, then $y=-x+3t$ for some integer $t$, and
$Q(x,y)\equiv3x^2\pmod9$. Thus $3\mid x,y$; the converse is immediate.
A component norm is divisible by $9$ exactly when that component is zero
modulo $3R$. Since $3\nmid n$, the other component is a unit of
$\overline R$. Its reduction is an Eisenstein unit, so the projective
reduction of the quaternion lies in $\Gamma_3$. The converse also holds.

Let $\delta$ be the indicator of the base point $\Gamma_3\in Y_3^+$,
and let $B_3(n)$ count the unfiltered representations with a component
norm divisible by $9$. Then
\[
B_3(n)=12(T_n\delta)(\Gamma_3).
\]
The projection of $\delta$ onto the constant and norm-sign characters
has diagonal value $1/3$, and the projection onto $W$ has diagonal
value $2/3$. Since $n\equiv1\pmod3$, both characters have eigenvalue
$\sigma(n)$. It follows that
\begin{equation}\label{eq:ramified-exact}
B_3(n)=4\sigma(n)+8\lambda_n.
\end{equation}

Return now to even $n=2^e m$. Let $B_3^o(n)$ count the same obstruction
among the retained odd-component representations. Scaling by $2$ preserves divisibility of either norm by
$9$, so $B_3^o(n)=B_3(n)-B_3(n/4)$, with the second term omitted
when $4\nmid n$. The unfiltered formula applies even when $n/4$ is odd.
If $n$ is not a square, both
$\lambda$ terms vanish. If $n$ is a square, then $e\ge2$ and
$\lambda_{n/4}=-\lambda_n/2$. Therefore
\[
B_3^o(n)=4S(n)+12\lambda_n,
\]
with the last term interpreted as zero for nonsquares. Dividing by
\eqref{eq:odd-count} gives
\begin{equation}\label{eq:ramified-odd-bound}
\frac{B_3^o(n)}{R_o(n)}
=\frac13+\frac{\lambda_n}{S(n)}
\le\frac13+\frac{2}{3\sqrt n}.
\end{equation}

\section{Putting numbers on the estimates}\label{sec:elementary}
The local estimate is in place. We now need bounds that can be added
over all unwanted primes while keeping the total below one. A bound on
norm multiplicities handles the large primes; elementary residue counts
handle how many small primes there are.

Define
\[
M(n)=\max_{1\le u\le n}\prod_{\substack{q^e\Vert u\\q\equiv1\ (3)}}(e+1).
\]
Unique factorization in $R$ shows that a positive integer $u$ has either
no Eisenstein norm representations or exactly
$6\prod_{q\equiv1\ (3)}(e_q+1)$ of them. Inert prime exponents must
be even, and the ramified prime contributes no choice. In particular
the count is at most $6M(n)$ for $u\le n$.

\begin{lemma}\label{lem:divisor}
For every positive integer $n$, one has $M(n)\le(15/8)n^{1/5}$.
\end{lemma}
\begin{proof}
The relevant local factors are $G_q(a)=(a+1)q^{-a/5}$ for split primes
$q\equiv1\pmod3$. For $q\ge32$, the inequality $a+1\le2^a$ gives
$G_q(a)\le1$. The successive ratio
$(a+2)/((a+1)q^{1/5})$ decreases with $a$, so each maximum is found by
checking up to the first $E$ with $(E+2)^5\le q(E+1)^5$.
For $q=7,13,19,31$, the maximizing exponents are $2,1,1,1$.
The product of the fifth powers of these maxima is
\[
\frac{7962624}{375193}<\left(\frac{15}{8}\right)^5.
\]
Multiplication over prime factors of each $u\le n$ proves the bound.
\end{proof}

\begin{lemma}\label{lem:lattice}
The number of nonzero Eisenstein integers of norm at most $t$ is at most
$6\lfloor t\rfloor\le6t$ for $t\ge0$.
\end{lemma}
\begin{proof}
Each orbit of the six units has a unique representative $x+y\omega$ with
$x\ge1$ and $0\le y<x$. Give this representative the label
$L=x(x-1)/2+y+1$. Labels are distinct positive integers. For $x\ge2$,
\[
L\le x(x+1)/2\le3x^2/4\le Q(x,y),
\]
and $x=1$ is immediate. Thus at most $\lfloor t\rfloor$ orbits occur.
\end{proof}
For an odd inert prime $p$, divisibility of a component norm by $p$
allows that component to be divided by $p$. Among the retained odd
components it is nonzero. Lemma~\ref{lem:lattice} gives at most $6n/p^2$
choices, and the other component has at most $6M(n)$ choices for each
norm value. Counting either spoiled component and dividing by
$R_o(n)=12S(n)\ge18n$ gives
\begin{equation}\label{eq:tail-reps}
\frac{b_p^o(n)}{R_o(n)}\le\frac{4M(n)}{p^2}
\le\frac{15}{2}\frac{n^{1/5}}{p^2}.
\end{equation}

Let $\PP_o=\PP\setminus\{2\}$. For $P\ge5$, we need
\begin{align}
\#\{p\in\PP_o:p\le P\}&\le\frac{2P}{15}+\frac{17}{15},\label{eq:pi}\\
\sum_{\substack{p\in\PP_o\\p>P}}p^{-2}
&\le\frac{2}{15P}+\frac{8}{5P^2},\label{eq:prime-tail}\\
\sum_{p\in\PP_o}\frac2{p(p-1)}&<\frac3{20}.\label{eq:prime-main}
\end{align}
Here are elementary proofs. Apart from $5$, these primes occupy the
residues $11,17,23,29\pmod{30}$. If $A(t)$ counts positive integers in
these four residues, the discrepancy $A(t)-2t/15$ is $30$-periodic, decreases linearly
between jumps at $11,17,23,29$, and has endpoint and left-limit bounds
\[
-\frac{22}{15}\le A(t)-\frac{2t}{15}\le\frac2{15}.
\]
This proves \eqref{eq:pi}, and also
$A(t)-A(P)\le2(t-P)/15+8/5$ for $t\ge P$. Partial summation gives the all-real bound
\[
\sum_{\substack{p\in\PP_o\\p>P}}p^{-2}
\le2\int_P^\infty\frac{A(t)-A(P)}{t^3}\,dt
\le\frac2{15P}+\frac8{5P^2}.
\]
Prime $5$ contributes the extra $1$ in the counting estimate and is
absent from this tail because $P\ge5$. Finally all odd inert primes are among $6k-1$,
$k\ge1$. Bounding the tail after $k=3$ by an integral yields
\[
\sum_{p\in\PP_o}\frac2{p(p-1)}
\le\frac1{10}+\frac1{55}+\frac1{136}
+\int_3^\infty\frac{2\,dx}{(6x-2)^2}
=\frac{6569}{44880}<\frac3{20}.
\]

\subsection{Leaving a decomposition unspoiled}\label{sec:infinite}
We can now put the two estimates together. Small primes use the spectral
bound, large primes use the lattice bound, and the mixed case adds the
condition at $3$. The aim is simply to make their combined fraction less
than one.

Assume \eqref{eq:reduced}, put $X=n^{1/4}$ and $P=2X$, and suppose
$X\ge100$. For the odd inert primes up to $P$, equations
\eqref{eq:odd-local}, \eqref{eq:pi} and \eqref{eq:prime-main} give
\[
\sum_{\substack{p\in\PP_o\\p\le P}}\frac{b_p^o(n)}{R_o(n)}
<\frac3{20}+\frac94 X^{-6/5}
\left(\frac{4X}{15}+\frac{17}{15}\right).
\]
For the larger primes, equations \eqref{eq:tail-reps} and
\eqref{eq:prime-tail} give
\[
\sum_{\substack{p\in\PP_o\\p>P}}\frac{b_p^o(n)}{R_o(n)}
\le\frac{15}{2}X^{4/5}
\left(\frac1{15X}+\frac2{5X^2}\right).
\]
The mixed case also excludes the ramified obstruction
\eqref{eq:ramified-odd-bound}. Its total spoiled fraction is therefore
less than
\begin{equation}\label{eq:final}
\frac{29}{60}+\frac{11}{10}X^{-1/5}
+\frac{111}{20}X^{-6/5}+\frac{2}{3X^2}.
\end{equation}
Since $(5/2)^5<100$, one has $X^{-1/5}<2/5$ and
$X^{-6/5}<1/250$, while $X^{-2}\le1/10000$. Thus \eqref{eq:final} is
strictly smaller than
\[
\frac{29}{60}+\frac{11}{10}\frac25+
\frac{111}{20}\frac1{250}+\frac2{30000}
=\frac{591}{625}<1.
\]
At least one odd-component representation survives.

If $n\equiv4\pmod6$, exactly one component norm is $0$ modulo $3$,
since an Eisenstein norm is $0$ or $1$ modulo $3$. Neither norm contains
an inert prime, and neither is divisible by $9$. The component divisible
by $3$ is therefore $3b$ with $b\in\AA$, and the other is in $\AA$.
This proves the mixed statement for reduced $n\ge10^8$.

If $n\equiv2\pmod6$, both component norms are $1$ modulo $3$.
The ramified obstruction is absent, so the same strict bound applies
after omitting its contribution. A surviving representation gives two
members of $\AA$. This proves the other binary statement in the same
infinite range. Both values admit primitive coordinate representations
by Proposition~\ref{prop:primitive-classification}; the surviving
quaternion coordinates themselves need not be primitive.

\section{Finishing the finite ranges}\label{sec:finite}
The estimates settle both reduced problems from $10^8$ onward. Below
that threshold we can check every target directly. The sieve is small
enough to reproduce, and its meaning is straightforward: each surviving
index records a positive integer in $\AA$.

For $c\in\{2,4\}$ put
\[
H_c=\left\lfloor\frac{10^8-c-1}{6}\right\rfloor.
\]
Index the possible summands by $6j+1$, $0\le j\le H_c$, and set
$\varepsilon_j=1$ unless a prime $p\equiv5\pmod6$, $p\le9999$, divides
$6j+1$. For such a divisor set $\varepsilon_j=0$; negative indices also
have value zero.

This is an exact membership test for $\AA$. Each indexed integer is odd
and prime to $3$. If it contains an inert prime factor, the total number
of inert factors counted with multiplicity is even, since the integer is
$1$ modulo $3$. At least one of these factors is at most the square root,
which is less than $10000$. For a prime $p\equiv5\pmod6$, precisely the
indices $(5p-1)/6+rp$, $r\ge0$, are removed.

\begin{proposition}\label{prop:finite-certificates}
For every target $n=6k+2<10^8$, there is an $i\le35$ with
$2i\le k$ and $\varepsilon_i\varepsilon_{k-i}=1$. For every target
$n=6k+4<10^8$, there is a $j\le33$ with $3j\le k$ and
$\varepsilon_j\varepsilon_{k-3j}=1$.
\end{proposition}
\begin{proof}[Computer-assisted proof]
The standard-library program printed in Appendix~\ref{app:verifier}
constructs the exact membership sieve, checks the small candidates by
direct divisor tests, and requires a witness at every target index.
Its two complete scans give the following results:
\begin{center}
\begin{tabular}{lrrr}
\toprule
family & targets & largest target & small summand bound\\\midrule
$\AA+\AA$ & $16{,}666{,}667$ & $99{,}999{,}998$ & $6i+1\le211$\\
$\AA+3\!\cdot\!\AA$ & $16{,}666{,}666$ & $99{,}999{,}994$ & $6j+1\le199$\\
\bottomrule
\end{tabular}
\end{center}
Neither scan has an uncovered target. The witnesses give, respectively,
\[
6k+2=(6i+1)+(6(k-i)+1),\qquad
6k+4=(6(k-3j)+1)+3(6j+1).
\]
All terms are positive members of the required sets. Separate NumPy
implementations regenerate the membership arrays and recheck every
witness. Their normalized membership and witness digests agree with the
printed verifier. The byte conventions and expected digests are given
in the appendix. The finite scans use integer arithmetic throughout.
\end{proof}
Proposition~\ref{prop:finite-certificates}, the infinite-range argument in
Section~\ref{sec:infinite}, and the initial multiplicative reduction prove
Theorem~\ref{thm:main}.

The scans find every witness using a small menu of summands. This raises
a further possibility: perhaps the same menu could work for every target.
The Chinese remainder theorem shows why these bounds have to remain
finite-range statements.

\begin{proposition}\label{prop:no-finite-menu}
For either $c=1$ or $c=3$, and every finite $F\subset\AA$, there are
infinitely many positive $n\equiv c+1\pmod6$ with
$n\notin\AA+(c\!\cdot\! F)$.
\end{proposition}
\begin{proof}
The case $F=\varnothing$ is immediate. Otherwise choose distinct odd
inert primes $q_f$, one for each $f\in F$. There are arbitrarily many
such primes: for any finite collection including $2$, the integer three
times their product minus one has an inert prime factor outside the
collection. The Chinese remainder theorem gives an arithmetic progression
of solutions to
\[
n\equiv c+1\pmod6,\qquad n\equiv cf\pmod{q_f}\quad(f\in F).
\]
Choose its positive terms larger than $c\max F$. Then $n-cf>0$ is
divisible by the forbidden prime $q_f$, so it is not in $\AA$, for every
$f\in F$.
\end{proof}
Thus the finite witness bounds $211$ and $199$ cannot be replaced by
fixed global bounds of the same kind.

\section{What the two binary answers give us}\label{sec:primitive-sumsets}
With both binary identities proved, the two families now fit into a
single familiar object. We look at norm values that admit coprime
coordinates:
\[
\NN=\{Q(x,y)>0:x,y\in\Z,\ \gcd(x,y)=1\}.
\]
For a set $S\subset\Z_{>0}$, let $kS$ denote the sum of exactly $k$
members of $S$, with repetitions allowed. We distinguish this from a
dilation, writing $3\!\cdot\! S=\{3s:s\in S\}$.

We begin by identifying these values, then use the two binary identities
to work out their sums.

\begin{proposition}\label{prop:primitive-classification}
The primitive Eisenstein norms are exactly
\[
\NN=\AA\sqcup(3\!\cdot\!\AA).
\]
\end{proposition}
\begin{proof}
If an inert prime divides $Q(x,y)$, it divides both coordinates by
\eqref{eq:anisotropic}. Also $9\mid Q(x,y)$ forces $3\mid x,y$.
Indeed, $3\mid Q(x,y)$ gives $x+y\equiv0\pmod3$; writing
$y=-x+3t$ gives
\[
Q(x,y)=3(x^2-3xt+3t^2).
\]
Divisibility by $9$ now forces $3\mid x$, and hence $3\mid y$.
Thus a primitive norm has no inert prime factor and contains at most
one factor $3$. It belongs to the stated union.

Conversely, $R=\Z[\omega]$ is norm Euclidean by the lattice rounding
argument, and hence is a unique factorization domain. For a rational
prime $q\equiv1\pmod3$, choose a nontrivial cube root
$t\in\F_q^\times$. The ideal $(q,\omega-t)$ has quotient $\F_q$;
its principal generator $\pi_q$ therefore has norm $q$. The two distinct
roots $t,t^{-1}$ yield conjugate, nonassociate factors. Thus
$q=\pi_q\bar\pi_q$, up to a unit.
For $a=\prod q^{e_q}\in\AA$, choose one factor $\pi_q$ at each $q$ and
put $z=\prod\pi_q^{e_q}$. Then $N(z)=a$. The element
$(1-\omega)z$ has norm $3a$. Neither element is divisible by a rational
prime: a split rational prime would require both conjugate factors,
an inert prime cannot divide its norm, and $3$ is associated to
$(1-\omega)^2$, whereas at most one factor $1-\omega$ was used.
Their coordinates in the integral basis $1,\omega$ are therefore
coprime. Finally, the two sets are disjoint, since elements of $\AA$
are prime to $3$.
\end{proof}

We can now use the two binary answers in their simplest form:
\begin{equation}\label{eq:binary-input-sumsets}
\AA+\AA=\{n>0:n\equiv2\pmod6\},\qquad
\AA+(3\!\cdot\!\AA)=\{n>0:n\equiv4\pmod6\}.
\end{equation}
The mixed identity fills the second family of residue classes. Together
with the first identity, it gives the complete two-summand answer below.

\begin{proposition}\label{prop:primitive-binary}
One has
\[
\NN+\NN
=\{n\in\Z_{>0}:n\text{ is even},\ n\not\equiv0,12\pmod{18}\}.
\]
\end{proposition}
\begin{proof}
Expand the sum using Proposition~\ref{prop:primitive-classification}.
The three possibilities and their exact positive residue classes are
\[
\begin{array}{c|c}
\text{summands}&\text{residue classes modulo }18\\ \hline
 a+b\quad(a,b\in\AA)&2,8,14\\
 a+3b\quad(a,b\in\AA)&4,10,16\\
 3a+3b\quad(a,b\in\AA)&6.
\end{array}
\]
The first two rows follow from \eqref{eq:binary-input-sumsets}; the last
is three times its first identity. Their union is the stated set.
\end{proof}
The missing classes also have a direct local explanation. Every primitive
norm is odd. If their sum is divisible by $3$, both summands belong to
$3\!\cdot\!\AA$; each is $3$ modulo $18$, so their sum is $6$ modulo
$18$. This excludes $0$ and $12$ modulo $18$.

Two summands are the hard part. For more summands, the presence of $1$
does most of the work: subtract some ones, use a binary decomposition,
and put the ones back.

\begin{proposition}\label{prop:higher-sumsets}
For every $k\ge2$, one has
\[
k\AA=\{n\in\Z:n\ge k,\ n\equiv k\pmod6\}.
\]
Consequently every positive integer is a sum of at most seven elements
of $\AA$, and seven is necessary. For every $k\ge3$, one has
\[
k\NN=\{n\in\Z:n\ge k,\ n\equiv k\pmod2\}.
\]
Every positive integer is then a sum of at most four primitive norms,
and four is necessary.
\end{proposition}
\begin{proof}
For $k\AA$, positivity and congruence are necessary. Conversely, subtract
$k-2$ copies of $1$ from an integer on the right. The remainder is at
least $2$ and congruent to $2$ modulo $6$, so it is a sum of two members
of $\AA$. For $n\ge7$, choose the unique $k\in\{2,3,4,5,6,7\}$ congruent
to $n$ modulo $6$; the formula applies. For $n<7$, use ones.
A sum of $r$ elements of $\AA$ equal to $25$ has $r\equiv1\pmod6$.
Since $25=5^2\notin\AA$, it cannot have fewer than seven terms, and
\[
25=13+7+1+1+1+1+1.
\]

Every primitive norm is odd. To prove the converse for three summands,
let $n\ge3$ be odd. If $n\equiv3\pmod6$, use the formula for $3\AA$.
If $n\equiv5\pmod6$, write $n=3+(a+b)$ using the first binary identity.
If $n\equiv1\pmod6$, then $n\ge7$ and the mixed identity writes
$n=3+(a+3b)$. All three summands in each case are primitive norms.
For $k>3$, subtract $k-3$ ones and apply the three-summand result.
Thus odd integers at least $3$ use three terms and even integers at least
$4$ use four; $1$ and $2$ use one and two ones. The integer $12$ is not
a sum of two primitive norms by Proposition~\ref{prop:primitive-binary},
and parity excludes one or three terms. Since $12=3+3+3+3$, four is sharp.
\end{proof}

The two sharp examples mark the end of this chain: the binary identities
give all the higher sumsets, and $12$ and $25$ show exactly where the
universal bounds stop.

\section{The part of the argument we can reuse}\label{sec:orbit-criterion}
The local sieves above average a bad event over an orbit, apply a spectral
bound, and add the local losses. This appendix isolates that argument,
allowing the inverse-unit mass needed when several ideal classes occur.
The arithmetic work identifies the orbit, counts the representations,
and leaves a strict margin below one.

The first statement isolates the averaging step. The second lists the
uniform bounds that turn it into an eventual, and potentially effective,
representation theorem.

\begin{proposition}\label{prop:orbit-averaging}
Let a finite group $G$ act on the right on a finite set $Y$ equipped
with a $G$-invariant probability measure $\pi$. Let $K\le G$, and let
$\Omega=y_0K$ be a $K$-orbit of positive measure. On $L^2(Y,\pi)$ put
$f=1_\Omega$ and $\mu=\pi(\Omega)$. Suppose $T$ is self-adjoint,
commutes with the $G$-action, and satisfies $T1=d1$ with $d>0$.
Let $E\subset L^2(Y,\pi)$ be a subspace such that
\[
\|Tv\|\le\Lambda\|v\|\quad(v\in E^\perp),\qquad
f-\mu1\in E^\perp,
\]
where $\Lambda\ge0$. Then
\[
\left|\frac{(Tf)(y_0)}d-\mu\right|
\le\frac\Lambda d(1-\mu).
\]
In particular, the same right side bounds the upper deviation.
\end{proposition}
\begin{proof}
Both $f$ and $Tf$ are right-$K$ invariant, so $Tf$ is constant on
$\Omega$ and
\[
(Tf)(y_0)=\frac{\langle f,Tf\rangle_\pi}{\mu}.
\]
Write $f=\mu1+h$. Then $h\perp1$ and
$\|h\|_\pi^2=\mu(1-\mu)$. Self-adjointness and $T1=d1$ remove the
cross terms, giving
\[
\langle f,Tf\rangle_\pi=d\mu^2+\langle h,Th\rangle_\pi.
\]
The spectral hypothesis bounds the absolute value of the last term by
$\Lambda\mu(1-\mu)$. Divide by $d\mu$.
\end{proof}

Global transitivity is unnecessary: only constancy on the orbit
$\Omega$ enters the averaging identity. In Section~\ref{sec:minus11},
the principal ideal-class fiber has probability $w=3/5$, and the inert
bad orbit has conditional density $\mu_p=2/(p(p-1))$ in that fiber.
Thus its measure in the proposition is $\mu=w\mu_p$.
The orthogonality to norm characters and the separate ramified
calculation are checked there explicitly.

\begin{proposition}\label{prop:arithmetic-criterion}
For each admissible positive integer $n\ge n_0$, let $\mathcal R_n$ be a finite
set of $R(n)>0$ candidate representations. Let $E_n$ be an exceptional
bad subset, and $B_{p,n}$ bad subsets indexed by a fixed set $\mathcal S$
of primes. Every representation outside
$E_n\cup\bigcup_{p\in\mathcal S}B_{p,n}$ is assumed to be good.
Suppose, uniformly in the admissible targets and primes, that
\[
\frac{|E_n|}{R(n)}\le\rho+\eta(n),\qquad
\frac{|B_{p,n}|}{R(n)}\le\mu_p+C_1n^{\delta-1/2},\qquad
\frac{|B_{p,n}|}{R(n)}\le C_2n^\delta p^{-2},
\]
where $C_1,C_2,\rho,\mu_p,\eta(n)\ge0$, $\delta<1/4$,
$\eta(n)\to0$, and
\[
\sum_{p\in\mathcal S}\mu_p\le M<\infty,\qquad \rho+M<1.
\]
Assume also, for every real $P\ge P_0>0$, that
\[
\#\{p\in\mathcal S:p\le P\}\le\alpha P+\beta,\qquad
\sum_{\substack{p\in\mathcal S\\p>P}}p^{-2}
\le\frac\gamma P+\frac\zeta{P^2},
\]
with nonnegative constants. For fixed $t>0$, put
\[
L_1=\alpha tC_1+\frac{\gamma C_2}{t},\qquad
L_2=\beta C_1+\frac{\zeta C_2}{t^2}.
\]
Every admissible $n\ge\max\{n_0,(P_0/t)^4\}$ satisfying
\[
\eta(n)+L_1n^{\delta-1/4}+L_2n^{\delta-1/2}<1-\rho-M
\]
has a good representation. In particular, every sufficiently large
admissible target does.

If the constants are effectively known and $\eta$ has an effective
modulus $N_\eta(u)$, meaning $\eta(n)\le u$ for all admissible
$n\ge N_\eta(u)$, put $\Delta=1-\rho-M$, $a=1/4-\delta$ and
$b=1/2-\delta$. An explicit sufficient cutoff is
\[
\max\left\{n_0,(P_0/t)^4,N_\eta(\Delta/4),
 (4L_1/\Delta)^{1/a},(4L_2/\Delta)^{1/b}\right\}.
\]
\end{proposition}
\begin{proof}
Use the first local bound for $p\le P=t n^{1/4}$ and the second for
$p>P$. The union bound gives
\begin{align*}
\frac{|E_n|+\sum_{p\in\mathcal S}|B_{p,n}|}{R(n)}
&\le\rho+M+\eta(n)+C_1n^{\delta-1/2}(\alpha P+\beta)
+C_2n^\delta\left(\frac\gamma P+\frac\zeta{P^2}\right)\\
&=\rho+M+\eta(n)+L_1n^{\delta-1/4}+L_2n^{\delta-1/2}.
\end{align*}
The tail bound makes the prime sum convergent. If the right side is
less than one, the bad subsets cannot cover $\mathcal R_n$.
Both powers tend to zero because $\delta<1/4$.
At the displayed effective cutoff, the three error terms are each at
most $\Delta/4$, giving a total at most $1-\Delta/4<1$.
\end{proof}
The prime-count hypotheses always admit the crude choices
$\alpha=\gamma=\zeta=1$, $\beta=0$, $P_0=1$, by comparison with integers
and a decreasing integral. Arithmetic residue information improves these
constants. When $\alpha C_1$ and $\gamma C_2$ are positive, the leading
coefficient $L_1$ is minimized at
\[
t=\sqrt{\frac{\gamma C_2}{\alpha C_1}}.
\]
This minimizes that coefficient, not necessarily the complete cutoff.
The mixed Eisenstein application uses $\delta=1/5$, $\rho=1/3$,
$\eta(n)=2/(3\sqrt n)$, $C_1=9/4$, $C_2=15/2$, $M=3/20$ and
$(\alpha,\beta,\gamma,\zeta)=(2/15,17/15,2/15,8/5)$.
Section~\ref{sec:infinite} uses $t=2$ and the sharper direct inequality
there to obtain $10^8$.

\section{The next question: balanced summands}\label{sec:balanced}
Existence still leaves the shape of a decomposition largely open. It is
natural to ask for the two summands to have comparable size.
The supplementary NumPy verifier for $\AA+\AA$ also tests $n/3\le a\le b\le2n/3$. In the entire range
$2\le n<100{,}000{,}000$, $n\equiv2\pmod6$, precisely $8$ and $44$
fail. This motivates, but does not prove, the sharper statement:
\begin{conjecture}
Every positive integer $n\equiv2\pmod6$ other than $8$ and $44$ admits $n=a+b$ with
$a,b\in\AA$ and $n/3\le a,b\le2n/3$.
\end{conjecture}
The present counting argument leaves the two norm sizes free. To settle
this last question, we would need control of the ratio
$N(z)/(N(z)+N(w))$, or a different construction that keeps the two values
in the desired interval. The finite search gives a clear target, but it
does not yet provide that extra argument.


\paragraph{AI assistance and review.}
ChatGPT (OpenAI) assisted with mathematical arguments, literature checks,
exposition, proof review and verification code.

\appendix
\section{A standalone verifier for both finite ranges}\label{app:verifier}
The program below uses only the Python standard library. Its complete
inputs are the two residue classes, the cutoff $10^8$ and the small
summand bounds in Proposition~\ref{prop:finite-certificates}. It constructs
the membership sieve itself. For every target it finds an allowed
complementary index, or raises an exception. The direct test of all
divisors of a small candidate is equivalent to requiring all its prime
factors to be $1$ modulo $3$: products of such primes have every divisor
$1$ modulo $3$.

Membership entries are ordered bytes $0$ or $1$. A witness index is
encoded as a signed little-endian 32-bit integer in increasing target
order. The witness is $i$ in the unrestricted identity and $j$, the index
of $b$, in the mixed identity. The expected SHA256 digests are:

\noindent\textit{Unrestricted membership:}
\begin{quote}\footnotesize\ttfamily
010a6f3752c25b6e58164d309c2173f4b308f5b81180ad96c3e96dfca0994575
\end{quote}
\noindent\textit{Unrestricted least-witness indices:}
\begin{quote}\footnotesize\ttfamily
fff9741cf68c12c795e28aebd94884b2717272dd337de5bd28f07234b658ddf0
\end{quote}
\noindent\textit{Mixed membership:}
\begin{quote}\footnotesize\ttfamily
aadd1a6ebe6f6450f4b7e270d308af6971f1b38d3b388b768c3d28c752518fc0
\end{quote}
\noindent\textit{Mixed least-$b$ indices:}
\begin{quote}\footnotesize\ttfamily
91496589fd9b1be7e9aba7cd77d63cb2c0de4be710564172ded09198c638c49d
\end{quote}
The membership arrays have different lengths because the two target
ranges have different final indices. The digests identify outputs;
correctness rests on the exact sieve and exhaustive witness tests.
The accompanying \href{source.zip}{\texttt{source.zip}} contains this
program as \texttt{verify\_binary\_certificates.py}, the separate NumPy
implementations, and the exact-constant and ramified-count checks.

\begingroup\footnotesize
\verbatiminput{verify_binary_certificates.py}
\endgroup


\begin{thebibliography}{9}\small
\bibitem{Euler} W. Howe, G. Huffman, J. Kohlmeyer, and L. Kruer,
\emph{Euler's 1747 conjecture on sums of integers with restricted prime
factors}, Hexagon 2611.00370v1, 12 November 2026.
\url{https://hexagonmath.org/2611.00370v1}.
\bibitem{EMV} J. S. Ellenberg, P. Michel, and A. Venkatesh,
\emph{Linnik's ergodic method and the distribution of integer points on
spheres}, in Automorphic Representations and $L$-Functions, TIFR Studies
in Mathematics 22 (2013), 119--185, Section 8.4.
\url{https://www.math.ias.edu/~akshay/research/gg.pdf}.
\bibitem{Li} W.-C. W. Li, \emph{Number theoretic constructions of Ramanujan
graphs}, Ast\'erisque 228 (1995), 101--120, Section 2.
\url{https://www.numdam.org/article/AST_1995__228__101_0.pdf}.
\bibitem{Voight} J. Voight, \emph{Quaternion Algebras}, Graduate Texts in
Mathematics 288, Springer, 2021, Chapters 28 and 41.
\url{https://doi.org/10.1007/978-3-030-56694-4}.
\bibitem{RV} O. Ramar\'e and G. K. Viswanadham,
\emph{The trigonometric polynomial on sums of two squares, an additive
problem and generalisation}, arXiv:2509.23260v1, 2025.
\url{https://arxiv.org/abs/2509.23260v1}.
\bibitem{JL} H. Jacquet and R. P. Langlands,
\emph{Automorphic Forms on $\mathrm{GL}(2)$}, Lecture Notes in Mathematics
114, Springer, 1970, Theorem~14.4 and Section~5.
\url{https://publications.ias.edu/sites/default/files/automorphic-forms-on-gl2_rpl_9.pdf}.
\bibitem{DeligneModular} P. Deligne,
\emph{Formes modulaires et repr\'esentations $\ell$-adiques},
S\'eminaire Bourbaki, expos\'e 355 (1968/69), published 1971, 139--172.
\url{https://www.numdam.org/item/SB_1968-1969__11__139_0.pdf}.
\bibitem{DeligneWeil} P. Deligne, \emph{La conjecture de Weil. I},
Publications Math\'ematiques de l'IH\'ES 43 (1974), 273--307.
\url{https://www.numdam.org/item/PMIHES_1974__43__273_0.pdf}.
\end{thebibliography}

\end{document}
